The corollary has a sharp empirical counterpart. At the apartment estimates of Section~\ref{ss:att_estimates}, taking the reform-period estimate $\hat\psi_\delta$ as a proxy for $\psi$ (exact under the maintained assumption of Proposition~\ref{prop:att_identification}), the band spans \attBandLoApt{} to \attBandHiApt{} log points, a width of \attBandWidthApt{}, against \attBandWidthBaseApt{} in the model of Section~\ref{sec:equilibrium}. The mean premium lies inside the aggregate band by construction of the aggregation, which is not evidence. Market by market, with a common $\psi$, the premium falls inside the market's own band in \attShareInBandApt{} percent of apartment markets: a single hiding cost per market would have to fall inside a window of about \attBandWidthApt{} log points for an apartment market to show interior hiding, yet the registry shows interior hiding rates in essentially every market. In houses the band is wider, \attBandWidthHouse{} log points, but only \attShareInBandHouse{} percent of house markets fall inside their own band. A single hiding cost per market therefore generates interior hiding only within a narrow window of costs, particularly in apartments; we do not model heterogeneity in hiding costs across sellers, and we read the band as a limitation of the single-cost specification.
